\section{Prior Work and Positioning}
\label{sec:related-work}
% -----------------------------------------------------------------------------
\subsection{Spectral theory, orthogonal invariants, and jet bundles}
The classification of the pair $(B,b)$ is not new as an abstract result. It follows from the spectral decomposition of a self-adjoint operator and the action of its orthogonal stabilizer on a distinguished vector. Spectral multiplicity theory, orthogonal matrix invariants, and vector invariants provide the classical framework \citep{halmos1951,weyl1946,sibirskii1967,sibirskii1968,khadzhiev1967}. Our contribution lies in selecting the conditional datum $(G_S,\Gamma_i,\alpha_i)$, interpreting its signature as an observation register, and explicitly factoring learning diagnostics through it. Likewise, the term ``jet'' is reserved for its standard meaning \citep{saunders1989}: the object studied here is a relative Gauss--Newton quadratic datum decorated by the context.

\subsection{Local influence, diagnostics, and robustness}
Leverage, studentized residuals, and Cook's distance belong to the tradition of regression diagnostics \citep{cook1977}. Differential geometry of perturbations and local influence was already developed by \citet{cook1986}, then connected to leverage and curvature in nonlinear regression by \citet{stlaurent1993}. Influence functions and infinitesimal robustness go back in particular to \citet{hampel1974}; Huber and bounded-influence regression estimators limit extreme residuals and leverage points \citep{huber1964,mallows1975,krasker1982,richardson1997}. The Versatile Influence Function now extends this analysis to non-decomposable objectives such as contrastive, ranking, or graph losses \citep{deng2025vif}; our current construction instead assumes an explicit decomposition into contributions of observations or declared blocks. Cellwise robust M-regression also addresses contaminated cells, a problem that norm clipping does not solve \citep{filzmoser2020}.

This paper therefore claims neither the leverage--residual combination nor the principle of bounded influence. Mallows and Schweppe forms belong to this tradition of GM-estimators \citep{mallows1975,richardson1997,maronna2006robust}. The candidate contribution is the orbit classification of the local triplet $(G,\Gamma,\alpha)$ and the simultaneous preservation of the observability operator and the effort before scalarization.

\subsection{Scores, Fisher geometry, tangent kernels, and representations}
The Fisher kernel of \citet{jaakkola1998} already constructs a similarity from cotangent scores in an information metric. Natural gradient and its parameterization-invariant geometries are prior work \citep{amari1998,ollivier2015,martens2020}; the relation between online natural gradient and the Kalman filter is made explicit by \citet{ollivier2017}. Degenerate pullback metrics, pseudometrics, and associated quotients for networks are studied directly by \citet{benfenati2023singular}; this prior work in particular delimits our main assumption $G_S\succ0$ and quotient-based extensions. NTK and GTK describe functional operators derived from the Jacobian and a parametric cometric \citep{jacot2018,bai2025gtk}. Evolution and alignment of tangent features are also studied by \citet{baratin2021}.

The register is therefore not the first geometric object built from scores or Jacobians. It differs through its conditioning on a single observation and through jointly retaining the positive second-order term $\Gamma_i$ and the first-order term $\alpha_i$, classified relative to the collective geometry.

\subsection{Data valuation and semivalues}
Data Shapley chooses a utility and averages marginal contributions over contexts \citep{ghorbani2019}. Beta--Shapley and Data Banzhaf modify the context distribution or the stability of this average \citep{kwon2022,wang2023banzhaf}. Volume-based methods explicitly address robustness to replication \citep{xu2021volume}; a recent preprint proposes geometric valuation by leverage and ridge leverage \citep{mendozasmith2025}. Datamodels and TRAK develop other approximations of the effect of data on predictions \citep{ilyas2022,park2023trak}.

Three works from 2026 further tighten this neighborhood. The Bayesian information approach of \citet{tailor2026bayesian} measures predictive information loss induced by data removal and approximates it using tangent features. Eigen-Value \citep{choi2026eigenvalue} values points through perturbations of covariance-eigenvalue ratios under domain shift. Finally, Symbolic Mechanistic Data Attribution \citep{habibi2026smda} fits Ridge regression on interpretable features and analytically decomposes data effects into an activation-mediated path and an output-mediated path. These works confirm that information, spectra, tangent features, and input--output separation are already active directions. Our result remains narrower: it jointly classifies the positive operator $B$ and the effort $b$, relative to context, before any scalar utility is selected.

The paper does not claim that Shapley must be scalar in principle: vector-valued utilities may be defined. Common uses of Data Shapley nevertheless choose a scalar utility before averaging and do not automatically retain the mixed tensorial structure $(\Gamma_i,\alpha_i)$. Our exact result is narrower: for additive information utility, the marginal contribution is a scalarization of the register; for a smooth predictive utility, the local effect or finite path provides the contribution before contextual averaging.

\subsection{Dynamic attribution}
Integrating a sensitivity along a participation path has direct precedent in dynamic valuation methods, notably NDDV \citep{liang2024nddv}. The Memory--Perturbation Equation likewise relates model memory during training to sensitivity under data perturbations and unifies several existing measures \citep{nickl2023mpe}. In its current formulation, NDDV connects finite marginal contribution to an integral of sensitivities along a state dynamics and discusses its relation to Shapley, Banzhaf, and leave-one-out. Our path theorem is an application of the implicit-function theorem and is not, in isolation, a deep novelty. The proposed distinction is that the path of minima is equipped at every instant with an observability--effort decomposition and a contraction operator. This distinction should be compared more directly with state--adjoint formulations and the MPE in future work.

\subsection{Boundary of the contribution}
\begin{statementbox}{Defensible claim}
The central result is neither a new robust estimator nor a new semivalue. It is the complete characterization, up to local reparameterization, of the relative quadratic model of an observation by the orbit signature of $(B,b)$, together with a proof that several common diagnostics are non-injective projections of this object. The path register extends this accounting to the case in which geometry and model evolve.
\end{statementbox}
Absolute bibliographic priority for a register with exactly this combination cannot be established from the search conducted here alone. The paper claims a theorem and a precise organization, not the general invention of geometric influence, dynamic valuation, or robustness.

% -----------------------------------------------------------------------------
